% ─── Abstract ─────────────────────────────────────────────────────────────────
This paper provides a merit-order assessment of green hydrogen's role in European
decarbonisation, across the EU-27, the United Kingdom and Switzerland, 29 countries over
2025 to 2050, built on a residential heat-demand reconstruction resolved to the third
level of the Nomenclature of Territorial Units for Statistics (NUTS3) and assessed at
the country level. The context is the buildings sector, which accounts for approximately
40 per cent of final energy consumption across the European Union and Switzerland and is
therefore central to the net-zero-by-2050 agenda
\citep{IEA2022buildings,EuropeanCommission2020GreenDeal}. Heat-pump electrification
leads the field \citep{IEA2022heatpumps,Abbasi2021}, yet the role of green hydrogen in
heating remains contested, and is usually settled by assuming a deployment share instead
of deriving one \citep{Rosenow2022PipeDream,Cebon2022}. We reconstructed residential
heat demand bottom-up from building footprints and archetype intensities and carried it
through a 200-draw Monte Carlo over four scenarios and a least-cost optimisation.

Hydrogen was priced as a system operator prices it, on short-run marginal cost in
dispatch, which bounds the heating share it could hold by where it wins the winter-peak
merit order. That bound is what the dispatch tests deliver; the share each scenario
realises is a policy target we set, and it did not move when the cost comparison did.
Three tests carry the assessment, cost competitiveness on a levelised basis, merit-order
dispatch, and investment viability or capital recovery (the missing-money test). The third
is decisive rather than incidental. No hydrogen peaker recovers its capital from market
rent in any country or scenario we run, at about 9 per cent capacity-weighted across the
markets that build. The paper then traces the electrified heat load into the power system and assesses
hydrogen as a dispatchable backup generator.

On the levelised-cost test the outcome settles, and it settles on the cost of holding
the molecule between seasons. Geology sets that cost for two of the seven surviving
hydrogen leads; the retail tariff on the rival carrier sets it for the other five. At 2050
medians of \euro{}118.5 for the best heat pump, \euro{}122.8 for the hydrogen boiler and
\euro{}132.7 for the gas boiler per MWh of useful heat, with each carrier charged its
own last-mile network and with hydrogen charged the seasonal store that its
winter-weighted load shape requires, the best heat pump is the cheaper of the two clean
options in 22 of 29 countries and the hydrogen boiler in seven, at a median country gap
of \euro{}15.4/MWh in the heat pump's favour (mean \euro{}19.7, and \euro{}1.4 weighted by each country's 2025 heat demand). Five of those seven are salt-cavern markets, including the four widest.
Space heating draws about three-quarters of its energy between October and March while
electrolytic production is far flatter, so the mismatch must be stored, and hydrogen's
base-load case concentrates where the geology makes storing it cheap.

That is the same condition that decides the peaking arena, so the base-load and dispatch
tests now agree on one physical quantity, the cost of holding the molecule between
seasons, which is also why their agreement corroborates less than two separate
confirmations would. The two assumptions still most generous to
hydrogen both run the same way. By 2050 the median study-area household pays roughly \euro{}209/MWh for
taxed retail electricity against roughly \euro{}50/MWh for hydrogen delivered and
untaxed, so equalising the fiscal treatment would widen the heat pump's margin, and
repricing the molecule as blue at the \euro{}81/MWh of our own blue route, which is the provenance the one dedicated assessment of heating hydrogen concludes it would have, adds a fuel wedge of \euro{}31/MWh, wider than the widest surviving hydrogen lead. The count turns on the hydrogen price
path and not the policy scenario, standing at 13, 22, 26 and 29 of 29 for the heat pump
as that path runs from rapid decline to stranded.

The dispatch test then bounds hydrogen across three merit-order arenas,
each pricing it on short-run marginal cost in a different market: the building
winter-peak heating stack, the district-heating stack, and the power-sector peaker
that supplies the electrified load. Under the most
favourable scenario (H2 Push, and all three counts below are on that scenario, charging
both carriers alike) hydrogen
wins the building winter-peak dispatch in 14 of 29 countries; undercuts gas
combined-heat-and-power in the district-heating stack in seven; and displaces the
carbon-priced gas peaker in the power sector that the electrified load runs on in seven
of 29, all of them salt-cavern markets, on the even-handed series that holds both
turbines at the same efficiency. Letting the two efficiencies diverge as each technology
improves, which hands hydrogen the better machine, raises that count to 15, but the
additional eight sit within \euro{}8/MWh-e of parity while the seven win decisively, so
seven is the count we carry. Where it wins, the dispatch role is secured by operating cost and
the carbon price on the fossil competitor, more than by technology or cheap hydrogen
alone; high achievement in one arena does not carry to another.

Even where hydrogen wins the dispatch competition, a capital-recovery hurdle remains, and
it is the binding one. The peaker runs too few hours to recover its capital from market
rent, the classic missing-money problem, so any deliberate build requires a capacity
payment. A technology-neutral payment does not deliver one. Across the five markets that build, the hydrogen peaker's standing payment exceeds the gas peaker's in four of them,
covering 96 per cent of the winning capacity, so a neutral auction procures unabated gas.
Selecting hydrogen instead requires conditioning the payment on emissions, worth about
\euro{}172/tCO$_2$ of displaced gas generation, or \euro{}239 if the gas rent is floored
at zero. Because electricity is priced on operating cost, hydrogen keeps its dispatch
role; what it lacks is the investment case. The power role is the clearest case. The firm peaking fleet
that the electrified heat load calls for is about 278~GW against a study-area system
peak of about 889~GW, roughly 31 per cent of peak capacity at about 139 full-load
hours a year, a capacity contribution. Those power-sector figures rest on one
realisation of a single synthetic weather year, which is the largest simplification
here, so they size a role and are not an adequacy assessment.

Infrastructure costing places this in
context. The cost-optimal decarbonisation pathway is led by district heat at about 53
per cent of 2050 supply and heat pumps at about 45 per cent, with residual fossil at
about 1.4 per cent; its hydrogen share of about a third of one per cent is bounded by a
ceiling that is zero in 27 of the 29 markets and so adds no independent evidence on
hydrogen. That pathway carries the carbon price of the second emissions trading
system for buildings and road transport (ETS2) inside the cost of every fossil option.
Tightening the emissions cap from a 75 to a 100 per cent reduction moves the
present-value system cost by 0.05 per cent, so the cost of the ambition is small.

The
peaking verdict is robust to the structural choices it once rested on. The power-sector result comes from the reduced-form dispatch and is cross-checked
against an explicit unit commitment, which reproduces it to within \euro{}1.1~bn with
its constraints off and \euro{}1.3~bn with them on, the near-zero 2050 retrofit is foresight-invariant under a myopic rolling-horizon test, and it
holds under an endogenous treatment of Switzerland with electricity-price feedback,
each of which strengthens the conclusion without changing it. The least-cost reading
is nonetheless a property of the cost-optimal frontier.

On the peaking and capital-recovery evidence, then, hydrogen's direct heating niche is
small. Its widest count, 20 of 29 in district heating, is a ceiling on a residential gas
tariff rather than a competitive score, and falls to 7 at a wholesale gas price, which is
what a district-heat operator would pay. The larger potential is indirect, as a candidate
clean peaker, but only under a payment conditioned on emissions. A neutral one buys gas.

\medskip
\noindent\textbf{Keywords:} hydrogen; residential heating; heat pumps;
decarbonisation; merit order; missing money; net zero; European Union;
Switzerland.

\smallskip
\noindent\textbf{JEL codes:} Q41, Q48, Q54, H23
